% ---------------------------------------------------------------------
% Parte 2 · TCP y UDP · Erwin · 2:00 (lo edita solo Erwin)
% ---------------------------------------------------------------------
TCP y UDP están en la capa de transporte. Los dos mueven datos entre procesos de distintas máquinas, pero con una idea muy diferente.
TCP no envía nada hasta abrir la conexión con un handshake de tres pasos, y después numera los segmentos para que todo llegue completo y en el orden correcto \cite[p.~7]{barghash2022}.
Eso tiene un precio: cada envío espera su confirmación, y si algo se pierde hay que mandarlo otra vez, así que se demora más.
UDP no hace handshake ni revisa si los datagramas llegaron \cite[pp.~7--8]{barghash2022}. Es más liviano y por eso se usa donde perder un poco de información no importa tanto.
Pasa en una videollamada de Meet: si se pierde un pedazo de imagen ya no sirve pedirlo de nuevo, porque ese momento ya pasó.
Para el panel necesitamos que el servidor también pueda hablar primero. WebSocket hace eso y va sobre TCP, y es lo que les explica Mario.
